\section{Positive perturbations on a fixed model}
\label{mmp:section}

The construction in this section does not use the numerical-dimension
hypothesis.  It provides one model on which small positive perturbations
of the canonical divisor have base locus of codimension at least three.

\begin{proposition}\label{mmp:approximation}
Let $X$ be a smooth projective variety of dimension $n\geq 3$ with
pseudo-effective $K_X$.  There are an effective ample $\mathbb Q$-divisor
$B$ on $X$, an integer $s\geq 0$, and a finite sequence of
$K$-negative divisorial contractions and $K$-flips
\[
 X\dashrightarrow Y
\]
with the following properties.
\begin{enumerate}[label=\textup{(\roman*)}]
\item $Y$ is projective, $\mathbb Q$-factorial, and terminal.  The map
$X\dashrightarrow Y$ extracts no divisors.
\item Put $L=K_Y$, let $B_Y$ be the transform of $B$, and set
\[
 t_j=2^{-j},\qquad M_j=L+t_jB_Y\quad (j>s).
\]
For each $j>s$ there is a closed subset $Z_j\subset Y$ of codimension
at least three such that $|kM_j|$ is defined by a Cartier divisor and
is base point free on $Y\setminus Z_j$ for every sufficiently divisible
positive integer $k$.
\item On a common smooth projective resolution
$p:W\to X$, $q:W\to Y$, compatible canonical divisors satisfy
\begin{equation}\label{mmp:comparison}
 p^*K_X=q^*K_Y+E,
\end{equation}
where $E\geq 0$ is $q$-exceptional and its support contains the strict
transform of every prime divisor of $X$ contracted by $X\dashrightarrow Y$.
\end{enumerate}
\end{proposition}

Here and below, ``sufficiently divisible'' allows the required divisor
index to depend on $j$.  No uniform Cartier index for an infinite
sequence of models will be needed.

\subsection{Discrepancies and codimension two}

Write $a(F,U)=a(F;U,0)$ with the normalization fixed in the
introduction.  We record the strictness statement needed for the
finiteness argument.

\begin{lemma}\label{mmp:strict}
Suppose that $U\dashrightarrow U^+$ is a $K$-negative divisorial
contraction or a $K$-flip between normal projective
$\mathbb Q$-Gorenstein varieties.  Write $h:U\to R$ for the
contraction and $h^+:U^+\to R$ for the other morphism, with $h^+$
the identity in the divisorial case.  For every prime divisor $F$
over these varieties,
\[
 a(F,U)\leq a(F,U^+).
\]
The inequality is strict if the center of $F$ on $R$ is contained
in the locus where either $h$ or $h^+$ has a positive-dimensional
fiber.
\end{lemma}

\begin{proof}
Take a common smooth projective resolution carrying $F$:
\[
\begin{tikzcd}[column sep=large,row sep=small]
 & W \arrow[dl,"p"'] \arrow[dr,"q"] & \\
 U \arrow[dr,"h"'] \arrow[rr,dashed] & & U^+ \arrow[dl,"h^+"] \\
 & R. &
\end{tikzcd}
\]
Write $g=h\circ p=h^+\circ q:W\to R$ and set
$D=p^*K_U-q^*K_{U^+}$.  This divisor is $q$-exceptional,
since the step extracts no divisors.  On every curve contracted by
$g$, the divisors $-p^*K_U$ and $q^*K_{U^+}$ have nonnegative
degree; in the divisorial case the second has degree zero.
Thus $-D$ is $g$-nef, and hence $q$-nef.  The negativity lemma
gives $D\geq 0$; see \cite[Lemmas~3.38--3.39]{KM98}.

To verify the support assertion, let $r\in R$ be
a closed point with a positive-dimensional fiber on either side.
There is a curve in $g^{-1}(r)$ mapping onto a curve in that fiber:
take a component dominating that curve and cut it by sufficiently
general ample divisors.  The relative ampleness signs give $D\cdot C<0$
for this lifted curve.  Consequently $g^{-1}(r)$ meets
$\operatorname{Supp}D$.

We use the fiber-support form of the negativity lemma
\cite[Lemma~3.39]{KM98}: the support of an effective divisor
antinef over a proper birational morphism to a normal variety either
contains or is disjoint from each fiber. In the present setting it can
also be seen directly on the reduced irreducible components of the
fiber. If such a component $T$ is not contained in $\operatorname{Supp}D$
but meets it, a Cartier multiple of $D$ restricts to a nonzero effective
Cartier divisor on the integral projective variety $T$. Its degree against
sufficiently many ample hyperplanes is positive, producing a curve
contracted by $g$ with positive $D$-degree. This contradicts antinefness.
This argument does not require $T$ to be normal or the fiber to be reduced:
only its support is at issue. Components of dimension zero already lie
in the support if they meet it. The fibers of $g$ are connected because
$R$ is normal and $g$ is proper and birational
\cite[Corollary~III.11.4]{Hartshorne77}. Containment therefore propagates
through their reduced irreducible components.

It follows that every fiber in question is contained in
$\operatorname{Supp}D$.  If the center of $F$ on $R$ is contained
in this locus, then $F\subset\operatorname{Supp}D$ on our chosen
resolution.  Finally,
\[
 a(F,U^+)-a(F,U)=\operatorname{coeff}_F D,
\]
which proves both assertions.
\end{proof}

\begin{lemma}[Stabilization in codimension two]\label{mmp:stabilization}
Let
\[
 X=U_0\dashrightarrow U_1\dashrightarrow U_2\dashrightarrow\cdots
\]
be a finite or infinite sequence of $K$-negative divisorial
contractions and $K$-flips of projective $\mathbb Q$-factorial
varieties, starting from a smooth $n$-fold, $n\geq 3$.
After finitely many steps, each remaining map is an isomorphism
outside closed subsets of codimension at least three on both sides.
\end{lemma}

\begin{proof}
The usual MMP properties give preservation of $\mathbb Q$-factoriality,
invariance of the Picard number under a flip, and a drop of one under
an elementary divisorial contraction
\cite[Propositions~3.36--3.37]{KM98}.
There are therefore only finitely many divisorial contractions.
Since the steps extract no divisors, only finitely many prime divisors
of $X$ can disappear; denote this finite set by $\mathcal D$.

All models are terminal.  This follows from
\cite[Corollaries~3.42--3.43]{KM98}, or directly from
Lemma~\ref{mmp:strict}: previously exceptional valuations retain
log discrepancy greater than one, and a newly contracted divisor
increases strictly from log discrepancy one.  Moreover, every
valuation exceptional over the smooth $X$ has integral log
discrepancy at least two.  Discrepancies never decrease along the
sequence.

\smallskip\noindent
\emph{The flipped loci.}
Consider a flip whose target has a codimension-two component $T$
of its flipped locus.  A terminal variety is smooth in codimension
two \cite[Corollary~5.18]{KM98}.  Blowing up $T$ at its generic
smooth point therefore defines a prime valuation $F$ with
\[
 a(F,U_{i+1})=2.
\]
Its center on the contraction base is contained in the locus of
positive-dimensional fibers of the flipped contraction.
Lemma~\ref{mmp:strict} implies
\begin{equation}\label{mmp:witness}
 a(F,X)\leq a(F,U_i)<2=a(F,U_{i+1}).
\end{equation}
Thus $F$ is not exceptional over $X$: it is an original prime divisor.
It is exceptional over $U_{i+1}$, so it belongs to $\mathcal D$.
After this occurrence its discrepancy is at least two forever,
and it cannot satisfy \eqref{mmp:witness} at a later step.
Each flip with such a target component therefore uses a distinct
member of a finite set.  After a finite prefix, all steps are flips
and every flipped locus has codimension at least three.

\smallskip\noindent
\emph{The flipping loci.}
We have removed codimension-two components on the target side. It remains
to remove them on the source side; only then can a general surface be
transported unchanged through every later finite stage. The argument
adapts the homological decrease used in the proof of
\cite[Lemma~3.1]{AHK07}, keeping only classes of algebraic cycles.
Put $d=n-2$ and
define the finite-dimensional real vector space
\[
 \mathcal A_d(V)
 =\operatorname{span}_{\mathbb R}
 \{[T]:T\subset V\text{ is a closed integral $d$-fold}\}
 \subset H_{2d}(V,\mathbb R).
\]
Cycle classes,
Borel--Moore localization, and their compatibility with restriction
to an open subset are as in \cite[Section~19.1]{Fulton98}.
The homology of the compact complex algebraic varieties here is
finite-dimensional and vanishes above their real dimension.

For a remaining flip $U\dashrightarrow U^+$ over $R$, remove the
images in $R$ of both exceptional loci and take inverse images.
This gives a common open subset $\mathcal U$.  Indeed, a proper
birational morphism to a normal variety is an isomorphism over its
quasi-finite locus.  Each exceptional locus consists of
positive-dimensional fibers; its image has dimension at most one
less than its own.  Since the source contraction is small and the
target exceptional locus has dimension at most $n-3$, the closed
complements $Z\subset U$ and $Z^+\subset U^+$ satisfy
\[
 \dim Z\leq d,\qquad \dim Z^+\leq d-1.
\]
The additional points removed outside the exceptional loci lie in
an isomorphism locus and have dimension at most $n-3$.

Restriction gives maps
\[
 \mathcal A_d(U)\longrightarrow
 H^{\mathrm{BM}}_{2d}(\mathcal U,\mathbb R)
 \longleftarrow\mathcal A_d(U^+)
\]
with the same image: intersect an integral $d$-fold with
$\mathcal U$ and take its closure on the other model, while a
$d$-fold contained in the complement restricts to zero.
The right-hand map is injective, by the localization sequence and
$H_{2d}(Z^+,\mathbb R)=0$.
Hence
\begin{equation}\label{mmp:rank}
 \dim\mathcal A_d(U^+)\leq\dim\mathcal A_d(U).
\end{equation}
If the flipping locus has a $d$-dimensional component $T$, its
class belongs to the kernel of the left-hand restriction and is
nonzero: for an ample Cartier divisor $A$,
\[
 \langle c_1(A)^d,[T]\rangle=A^d\cdot T>0.
\]
Then \eqref{mmp:rank} is strict.  A nonnegative integer can drop only
finitely many times, so eventually the flipping loci also have
codimension at least three.  The same construction of the common
open then gives this codimension bound on both complements.
\end{proof}

\subsection{Construction of the perturbations}

\begin{proof}[Proof of Proposition~\ref{mmp:approximation}]
Choose an effective ample rational divisor $B$ for which $(X,B)$
is klt and $K_X+B$ is ample.  For example, take a sufficiently
positive general smooth very ample divisor and multiply it by a
rational number strictly between zero and one.  Since $K_X$ is
pseudo-effective, $K_X+tB$ is big for every rational $t>0$.

We shall repeatedly use a simple consequence of normality.  If a
birational map $X\dashrightarrow U$ extracts no divisors, every
prime divisor of $U$ corresponds to a prime divisor of $X$.
For compatible transformed divisors, the divisorial criterion for
regularity of a rational section consequently gives
\begin{equation}\label{mmp:sections}
 H^0(X,kD)\hookrightarrow H^0(U,kD_U)
\end{equation}
whenever the multiples are Cartier.  Thus effectiveness and bigness
of $B_U$, and bigness of $K_U+tB_U$, persist on every model
constructed below.  In particular, we do not need $B_U$ to be ample.

\smallskip\noindent
\emph{Finite stages.}
Set $t_j=2^{-j}$ for every integer $j\geq0$.
Starting with $V_0=X$, construct $V_{j+1}$ from $V_j$ by running
the MMP for
\[
 K_{V_j}+\Delta_j,\qquad
 \Delta_j=t_{j+1}B_{V_j},
\]
with scaling of
$C_j=(t_j-t_{j+1})B_{V_j}$.
The precise input is \cite[Corollary~1.4.2]{BCHM10}:
for a projective $\mathbb Q$-factorial klt pair $(V,\Delta)$
with big boundary $\Delta$, and $C\geq0$ such that
$(V,\Delta+C)$ is klt and $K_V+\Delta+C$ is nef,
the MMP with scaling of $C$ terminates.
Here the inductive hypothesis is that $(V_j,t_jB_{V_j})$ is
klt and its adjoint is nef.  The smaller pair is klt because
$B_{V_j}$ is effective; its boundary is big by
\eqref{mmp:sections}; and the larger pair is exactly the one in
the inductive hypothesis.

The adjoint remains big by \eqref{mmp:sections}, so a Mori fiber
space cannot occur: an effective representative of a positive
multiple cannot have negative degree on curves covering general
fibers.  Each stage is consequently a finite, possibly empty,
sequence ending at a nef adjoint.  Standard MMP preservation gives
the next klt pair and a projective $\mathbb Q$-factorial model.
The rational divisor $K_{V_j}+t_jB_{V_j}$ is nef and big, hence
semiample by the klt base point free theorem
\cite[Theorem~3.3]{KM98}.

\smallskip\noindent
\emph{Canonical steps.}
To apply Lemma~\ref{mmp:stabilization}, we must identify every
adjoint step as a canonical MMP step.  This requires canonical
negativity on the contracted ray and, for a flip, canonical
ampleness on the flipped side.
On its source $U$, let $D=K_U+t_{j+1}B_U$.  The contracted ray
$R_0$ satisfies
\[
 D\cdot R_0<0,\qquad
 (D+\lambda C)\cdot R_0=0,\qquad \lambda>0,
\]
where $C=(t_j-t_{j+1})B_U$ is the current scaling divisor.
Therefore $B_U\cdot R_0>0$ and $K_U\cdot R_0<0$.
This is the usual ray calculation for a directed program
\cite[proof of Corollary~1.3.3]{BCHM10}, and it already identifies
each divisorial step as a $K$-negative divisorial contraction.

For a small contraction $h:U\to R$, the adjoint flip makes $D^+$
relatively ample.  We show that $K_{U^+}$ is relatively ample as
well.  Choose the positive rational
number
\[
 b=\frac{K_U\cdot R_0}{D\cdot R_0}.
\]
The divisor $K_U-bD$ is numerically trivial over $R$ and descends
up to rational linear equivalence.  To see this, a Cartier
multiple is relatively semiample by the relative klt base point
free theorem for the pair $(U,t_{j+1}B_U)$
\cite[Theorem~3.24]{KM98}: subtracting its adjoint $D$ gives
relative ampleness, because $-D$ is relatively ample.  Its
semiample morphism is constant on each connected projective fiber,
as its degree on every fiber curve is zero.  The resulting image
is proper and quasi-finite over $R$, because each connected fiber
maps to one point. It is therefore finite and birational over the
normal $R$, hence equals $R$.
Thus $K_U-bD\sim_{\mathbb Q}h^*A$ for a rational Cartier divisor
$A$ on $R$.

Both morphisms in the adjoint flip are small, so this relation
transforms to
\[
 K_{U^+}-bD^+\sim_{\mathbb Q}(h^+)^*A.
\]
Since $D^+$ is relatively ample and $b>0$, so is $K_{U^+}$.
The adjoint flip is therefore a $K$-flip.  Thus all steps of the
finite stages satisfy the hypotheses of
Lemma~\ref{mmp:stabilization}.

\smallskip\noindent
\emph{The fixed model.}
Concatenate the finite stages and apply
Lemma~\ref{mmp:stabilization}.  Choose an endpoint $Y=V_s$ after
the resulting finite prefix.  If only finitely many nonempty stages
occur, choose $s$ after the last one.  For every $j>s$, the finite
composition $Y\dashrightarrow V_j$ is an isomorphism outside
codimension-at-least-three closed subsets on both models.
This property is preserved under finite composition: on a common
open, remove the next bad set and take its closure on the preceding
model; its dimension does not increase.  Let $Z_j$ be the resulting
closed subset of $Y$.

The divisor $M_j=K_Y+t_jB_Y$ agrees on this common open with
$K_{V_j}+t_jB_{V_j}$.  Choose a multiple Cartier on both varieties
and base point free on $V_j$.  A line bundle on a normal variety
has the same global sections after deleting a subset of codimension
at least two, since it is reflexive.  The two spaces of sections
are therefore identified, and the complete system on $Y$ is
generated outside $Z_j$.  This proves \textup{(ii)} for each $j$.

Finally, $Y$ is projective, $\mathbb Q$-factorial and terminal by
the preceding construction, and a finite composition of these
steps extracts no divisors.  On a common smooth projective
resolution of the finite chain, the effective differences in
Lemma~\ref{mmp:strict} telescope to \eqref{mmp:comparison}.
A divisor mapping onto a prime divisor of $Y$ also corresponds to
a prime divisor of $X$, so its coefficient in $E$ is zero.
Thus $E$ is $q$-exceptional.  For an original prime divisor
contracted on $Y$, terminality gives log discrepancy greater than
one on $Y$, compared with one on $X$.  Its strict transform has
positive coefficient in $E$, proving \textup{(iii)}.
\end{proof}
